% SPDX-License-Identifier: CC-BY-4.0
% Copyright (c) 2026 Martin Schröder <info@swedishembedded.com>

\section{The headline experiment and its replication: registered, not yet run}\label{sec:pending}

\paragraph{The claim this paper exists to establish.} Verified engineering-agent experience
can produce a sustained improvement in completion of previously unseen repository tasks,
under a controlled compute budget and without unacceptable loss of earlier capability.\cl{PR09}
\textbf{That claim is not yet measured.} No coding-task experiment has been run; this
section states exactly what will be run, what each outcome will mean, and where its numbers
will appear. Nothing in it is a result.

The experiments are registered in \code{prereg/coding-tasks.md}, dated 2026-10-16 and written
before any of their runs. \Cref{tab:registered} lists them. Their
numbers will be read from \code{data/results.csv} by \code{scripts/build_data.py}; a row exists
only for a finished run, and the status column below changes by itself when one does.\cl{PL02}

\begin{table}[t]
\centering\small
\input{generated/registered_table}
\caption{The registered coding-task experiments. The status is derived from the rows of
\code{data/results.csv}.}\label{tab:registered}
\end{table}

\begin{description}
\item[CE1, the untrained agent.] The base agent's resolve rate on repository tasks, its failure
profile, its cost per resolved task and the discordance between two seeds of the same arm per
family, which sizes the qualification partition. No threshold; nothing is trained before it is
reported.\cl{PR10}
\item[CE4, the headline.] A LoRA adapter on Qwen3-8B trained only on verified trajectories from
development-partition families, compared with the base on the qualification partition by the
primary test of \cref{sec:protocol}, served through Brain and driven by Sven, with retention,
serve-verdict agreement and the three controls. The headline holds only if the sign test has
$p<0.05$, retention is not refused and agreement is within the serve check's tolerance; any
other outcome is reported as what it is, including ``no held-out improvement'' and
``memorisation''.\cl{PR11}
\item[CE6, the 27B model.] Ten stages from loading the exact model to comparing correctness,
tool-call validity and resources against the base, gated by training correctness and by
deployment agreement. The 27B is not used for a confirmatory query until both gates pass;
the training rate of \cref{sec:throughput} is reported as measured and no target is
registered.\cl{PR12}
\item[CE7, replication.] The unchanged CE1 and CE4 procedure on a second, independent
repository chosen by a rule fixed in advance (permissive licence, offline build, a suite of
bounded running time, no part of the pipeline tuned on it), with the cross cells that separate
learning a repository from learning the protocol. The headline counts as reproduced only if
both repositories pass the primary test independently.\cl{PR13}
\end{description}

\begin{table}[t]
\centering\small
\input{generated/results_table}
\caption{Paired outcomes over qualification-partition families (wins: the adapter resolves more
of the family than the base; losses: fewer), with the exact one-sided sign test. Every cell is
empty until its run exists; a cell is never filled with an estimate.}\label{tab:results}
\end{table}

\paragraph{What a failure would show.} An adapter that does not win, a gain that a different
repository's adapter or unverified trajectories match, a retention loss beyond the margin, or
a served adapter that disagrees with the in-process one are each stated outcomes of the
registration, and each is publishable as such. A development gain with no qualification gain
(sign test $p>0.2$) is reported as memorisation. If fewer than 25 non-tied families are
available the comparison is labelled a pilot and the rule is not loosened.\cl{PR08}
